\begin{initiativkarte}{Bunte Hände}

  % Bild (optional)
  \IfFileExists{bilder/bunte_haende.jpg}{%
    \begin{center}
      \includegraphics[width=\linewidth,height=5cm,keepaspectratio]{bilder/bunte_haende.jpg}
    \end{center}
    \vspace{0.4cm}
  }{}

  % Kurzbeschreibung
  \textit{\textcolor{hawgray}{Bunte Hände ist eine studentische Initiative der HAW Hamburg, die sich für geflüchtete, internationale und deutsche Studierende sowie für eine offene Gesellschaft engagiert.}}

  \vspace{0.4cm}
  \hawrule

  % Beschreibung
  \textbf{\textcolor{hawblue}{Was sind die Bunten Hände?}}\\[0.3cm]

  Die Initiative \textbf{Bunte Hände} wurde 2018 an der HAW Hamburg gegründet und wird von der Arbeitsstelle Migration begleitet.

  Sie vereint geflüchtete, internationale und deutsche Studierende, die sich gemeinsam für Bildungsgerechtigkeit, Integration und gesellschaftliche Offenheit einsetzen.

  \vspace{0.3cm}

  \textbf{Unsere Schwerpunkte:}

  \begin{itemize}
    \item Unterstützung geflüchteter und internationaler Studierender beim Studienstart
    \item Freizeit-, Kultur- und Gemeinschaftsangebote
    \item Politisches Engagement für eine offene, demokratische Gesellschaft
    \item Fachliche Unterstützung im Studium
  \end{itemize}

  \vspace{0.4cm}

  \textbf{Campus Global – Projekte}

Ein zentrales Angebot ist mit verschiedenen Unterstützungsformaten: Study Guides bieten mehrsprachige Studienberatung (z. B. Arabisch, Kurdisch, Persisch), die Schreibwerkstatt unterstützt bei wissenschaftlichen Texten, Bewerbungen und Stipendien, und die Informatik-Werkstatt hilft bei Programmiersprachen wie C, C++, Java und Ruby sowie bei technischen Grundlagen.
  
  \vspace{0.2cm}

  \textbf{Freizeit \& Community}

  Neben der Studienunterstützung organisiert Bunte Hände regelmäßig:

  \begin{itemize}
    \item Spiele- und Kinoabende
    \item Semesterauftakt- und Abschlussfeiern
    \item Exkursionen und Events
    \item gemeinsames Iftar im Ramadan
  \end{itemize}

  \vspace{0.2cm}

  % Tags
  \tag{Diversität}
  \tag{Integration}
  \tag{Studienhilfe}
  \tag{Community}
  \tag{Politik}

  \vspace{0.2cm}
  \hawrule

  % Infozeilen
  \infozeile{\faUsers}{Zielgruppe: Geflüchtete, internationale und deutsche Studierende}
  \infozeile{\faGlobe}{Website: \href{https://www.bunte-haende.de/}{bunte-haende.de}}

\end{initiativkarte}
